\begin{titlepage}
    \begin{center}
        \heiti \bfseries \zihao{-1} \titleZh
        \\
        \vspace*{36pt}
        \heiti \zihao{-3}
        \begin{tabular}{ll}
            \makebox[7em][s]{专 \hfill 业~：} &
            \makebox[10em][l]{\MAJOR} \\
            \makebox[7em][s]{博 \hfill 士 \hfill 研 \hfill 究 \hfill 生~：} &
            \makebox[10em][l]{\STUDENT} \\
            \makebox[7em][s]{指 \hfill 导 \hfill 老 \hfill 师~：} &
            \makebox[10em][l]{\TEACHER~教授} \\
        \end{tabular}
        \vspace*{36pt}
    \end{center}

    \songti \zihao{-4}

    \section*{中文摘要}
    \phantomsection
    \addcontentsline{toc}{section}{中文摘要}

    \subsection*{背景：}
    在传统媒体向平台化传播迁移的过程中，东亚女性艺人的职业路径发生了显著变化。原本由杂志、影视和综艺主导的曝光体系，正在被社交媒体与视频平台重新组织。如何在多平台环境中识别职业阶段、解释危机后的形象恢复，并进一步构建可解释的转型预测框架，成为公开案例研究的重要问题。

    \subsection*{目的：}
    \begin{enumerate}[label=\textbf{\arabic*.}]
        \item 建立可复用的跨媒介职业生命周期模型；
        \item 总结影响职业转型成功与失败的关键机制；
        \item 以公开案例为基础，演示一个可解释、可扩展的预测框架写法。
    \end{enumerate}

    \subsection*{方法：}
    本研究使用公开报道、平台时间线、节目档期与访谈资料构建案例数据库，并以“阶段划分-事件编码-能力映射-韧性评分”的四层流程完成分析。具体方法包括：建立职业时间线、比较平台能力矩阵、提炼危机事件后的修复策略，以及使用规则驱动的评分函数生成示范性预测结果。

    \subsection*{结果：}
    研究发现，案例的职业转型具有明显的阶段性：早期依赖传统媒体建立识别度，中期通过影视与综艺实现角色扩张，危机期面临叙事权流失，重启期则依赖自有平台重建内容节奏。进一步分析表明，平台自主性、角色多样性与持续表达能力是影响转型稳定性的三项核心因素。基于这三类因素构建的示例预测框架能够较好地区分“稳定扩张”“短期波动”和“危机后重启”三种路径。

    \subsection*{结论：}
    跨媒介职业生命周期模型有助于把复杂案例拆解为若干可比较阶段，而可解释预测框架则为后续研究者提供了从描述走向预测的写作接口。对模板项目而言，这一公开示例证明了无需真实论文正文，也可以完整展示博士论文样式的封面、摘要、目录、章节、图表、附录与参考文献结构。

    \subsection*{关键词：}
    \begin{flushleft}
        \hspace{1em}\songti{\zihao{-4}{跨媒介转型，职业生命周期，形象韧性，公开案例研究，可解释预测框架}}
    \end{flushleft}

\clearpage{\cleardoublepage}
\end{titlepage}
